\documentclass[12pt]{article}
\usepackage[T1]{fontenc}
\usepackage[utf8]{inputenc}
\usepackage{lmodern}
\usepackage{textcomp}
\usepackage{enumitem}
\usepackage{geometry}
\geometry{margin=1in}
\begin{document}
\begin{center}
Answer Key - Health Sciences - Undergraduate Level Assessment 8 \
\small (Answers are ordered exactly as in the question paper.)
\end{center}

\section*{Multiple Choice}
\begin{enumerate}[label=\arabic*.]
\item \textbf{Answer:} D
\\ \textit{Options:} A) Orthodox Jews; B) Fundamentalist Protestants; C) Eastern Orthodox Christians; D) all of the above
\\ \textit{Note:} The position that abortion is murder is held by adherents of various religious traditions, including Christianity, Islam, and Judaism, as they often view the sanctity of life as a fundamental principle, leading to a shared belief that abortion equates to the taking of a life.
\item \textbf{Answer:} B
\\ \textit{Options:} A) Adipose; B) Cartilage; C) Epithelial; D) Muscle
\\ \textit{Note:} Cartilage is a type of flexible connective tissue that serves to cushion and support joints by providing a smooth surface for movement between bones.
\item \textbf{Answer:} C
\\ \textit{Options:} A) 30 through 50; B) at birth; C) 10 through 20; D) 20 through 30
\\ \textit{Note:} Fetishism typically develops during the formative years of adolescence, between the ages of 10 and 20, as individuals begin to explore their sexual identities and form attachments to specific objects or stimuli.
\item \textbf{Answer:} B
\\ \textit{Options:} A) Semantic memory; B) Knowledge about memory; C) Long-term memory for major events; D) All memory components together
\\ \textit{Note:} Metamemory is defined as an individual's awareness and understanding of their own memory processes, which encompasses knowledge about how memory works, including memory strategies, strengths, and weaknesses.
\item \textbf{Answer:} A
\\ \textit{Options:} A) Alpha-thalassaemia; B) Beta-thalassaemia; C) Hereditary persistence of fetal haemoglobin; D) Sickle cell disease
\\ \textit{Note:} Severe anaemia at birth is a feature of alpha-thalassaemia because this genetic disorder results in reduced production of alpha-globin chains, leading to ineffective erythropoiesis and significant haemolytic anaemia in affected newborns.
\end{enumerate}

\section*{True / False}
\begin{enumerate}[label=\arabic*.]
\setcounter{enumi}{5}
\item \textbf{Answer:} False
\\ \textit{Options:} A) True; B) False
\\ \textit{Note:} The statement is false because institutional care often emphasizes support and assistance rather than fostering independence, leading to a reliance on staff for daily activities rather than promoting self-sufficiency among older residents.
\item \textbf{Answer:} True
\\ \textit{Options:} A) True; B) False
\\ \textit{Note:} Asphyxiophilia involves the practice of asphyxiation for sexual arousal, which can lead to accidental death or serious injury due to the restriction of oxygen, making it one of the most dangerous sexual variations.
\item \textbf{Answer:} True
\\ \textit{Options:} A) True; B) False
\\ \textit{Note:} The inability to vaccinate in some countries has significantly hindered global polio eradication efforts because it creates pockets of unvaccinated populations that allow the virus to persist and spread, undermining worldwide immunization campaigns.
\item \textbf{Answer:} True
\\ \textit{Options:} A) True; B) False
\\ \textit{Note:} The statement is true because both genetic predispositions (nature) and environmental influences (nurture) interact to shape an individual's physical, emotional, and cognitive development throughout their life.
\item \textbf{Answer:} True
\\ \textit{Options:} A) True; B) False
\\ \textit{Note:} Providing comfort and reminiscing about happier times are both forms of emotional social support because they help individuals cope with stress and enhance their emotional well-being through positive interactions and shared memories.
\end{enumerate}

\section*{Long Form Response}
\begin{enumerate}[label=\arabic*.]
\setcounter{enumi}{10}
\item \textbf{Answer:} The high risk of HIV transmission from an infected mother to her fetus or infant in Africa is influenced by a combination of socio-economic, healthcare, and cultural factors. Socio- economically, poverty and lack of access to education can limit a mother's knowledge of HIV transmission and prevention methods. Healthcare systems in many African countries are often under-resourced, leading to inadequate prenatal care and antiretroviral treatments, which are critical in reducing transmission rates. Culturally, stigma surrounding HIV can discourage women from seeking testing and treatment, perpetuating cycles of transmission. Together, these factors create an environment where the risk of mother-to-child transmission remains alarmingly high.
\\ \textit{Note:} The answer correctly identifies that the interplay of socio-economic barriers, insufficient healthcare resources, and cultural beliefs significantly heightens the risk of HIV transmission from mothers to their infants in Africa, emphasizing the need for comprehensive interventions in these areas.
\item \textbf{Answer:} Saliva plays a crucial role in the digestion of starches, primarily through the action of the enzyme amylase, specifically salivary amylase or ptyalin. When food is chewed, saliva mixes with it, and salivary amylase begins to break down starch molecules into simpler sugars, such as maltose. This enzymatic action occurs in the mouth and initiates the carbohydrate digestion process before the food reaches the stomach. The mechanism of action involves the hydrolysis of the glycosidic bonds in starch, effectively converting complex carbohydrates into more easily absorbable forms. Thus, saliva not only facilitates the mechanical breakdown of food but also initiates the biochemical process of starch digestion.
\\ \textit{Note:} correct because it accurately describes the role of saliva and the specific enzyme salivary amylase in the initial digestion of starches by hydrolyzing glycosidic bonds to convert starches into simpler sugars in the mouth.
\end{enumerate}

\vfill
\noindent\textit{End of Answer Key}
\end{document}